\chapter{結論}
\chapoverview{本章では，本研究の目的，手法，結果，貢献，限界をまとめる．本研究は，特許文書を課題文脈と解決手段文脈に分離し，課題が近い異分野間において，対象分野では低出現である技術手法を専門家確認候補として抽出する枠組みを検討した．最後に，卒論段階で達成できた点と，今後の研究課題を整理する．}

\section{本研究のまとめ}
本研究の目的は，特許文書から課題文脈と解決手段文脈を分離して抽出し，課題が近い異分野間において，対象分野では低出現または未使用である他分野の技術手法を，専門家が確認すべきクロスドメイン技術候補として抽出できるかを検討することであった．

本研究では，医療画像解析を参照分野，製造欠陥検出を対象分野とし，U-Netを主ケースとして扱った．2015--2018年を過去期間，2019--2024年を将来評価期間とし，過去期間の手法ギャップから候補を抽出し，後年の対象分野で出現したかを確認する後ろ向き評価を設計した．

\section{主な結果}
本研究で得られた主な結果は以下の3点である．

第1に，全期間において，医療画像解析分野のU-Net出現率は製造欠陥検出分野より高かった．Publication numberベースでは，医療画像解析分野で5.30\%，製造欠陥検出分野で0.59\%であった．Family IDベースでも，医療画像解析分野で6.78\%，製造欠陥検出分野で0.55\%であった．

第2に，製造欠陥検出分野では，2015--2018年にU-Net関連特許が確認されず，2019--2024年に72件確認された．Family IDベースでも，2015--2018年は0件，2019--2024年は49件であった．

第3に，代表特許の目視確認により，医療画像解析側では腫瘍や生物医学画像のセグメンテーションに，製造欠陥検出側では管路，鋼材，金属表面欠陥のセグメンテーションまたは検出に，U-Net系構造が技術的手段として用いられていることを確認した．

\section{本研究の貢献}
本研究の貢献は以下の3点である．

第1に，特許文書を課題文脈と解決手段文脈に分離し，クロスドメイン技術候補抽出を，全文類似ではなく，課題の近さと解決手段の低出現性として整理した点である．

第2に，U-Netを対象に，医療画像解析と製造欠陥検出の特許データを比較し，手法ギャップと将来出現を用いた後ろ向き評価の設計を示した点である．

第3に，Publication numberベースとFamily IDベースの両方で傾向を確認し，代表特許の目視確認を組み合わせることで，候補抽出を単なる語の一致に留めず，専門家確認候補として解釈可能かを検討した点である．

\section{限界と今後の展望}
本研究は，技術導入可能性を直接証明するものではない．特許データ上での出現率差やBERT類似度は，あくまで専門家確認候補を絞り込むためのスクリーニング指標である．実際の導入可能性，性能，コスト，権利範囲，事業化可能性は，別途専門家が評価する必要がある．

今後は，YOLO，GAN，Domain Adaptationなどの補助ケースを追加し，提案枠組みの一般性を検討する必要がある．また，花王の洗浄・除去系技術から半導体洗浄への展開のような発明主題レベルの異質解決策型分析を行うことで，単一手法名に依存しない候補抽出へ拡張できる可能性がある．さらに，論文データと特許データを接続することで，科学知識が産業応用へ展開する経路を分析することも今後の重要な課題である．

\section{最終的な結論}
本研究は，U-Netの医療画像解析から製造欠陥検出への展開を対象に，特許文書を課題文脈と解決手段文脈に分離して扱うクロスドメイン技術候補抽出枠組みを検討した．予備分析の結果，U-Netは医療画像解析分野で先行して出現し，製造欠陥検出分野では後年に出現・増加していた．この結果は，提案枠組みが，既知のクロスドメイン応用事例を専門家確認候補として抽出できる可能性を示すものである．
